%% ma: L11-B15-0093
%% lop: 11
%% chuong: 5
%% bai: 15
%% dang: 11.15.5
%% loai: TLN
%% muc: VDC
%% dap_an: "12"
%% nguon: "tai-lieu-hoc-tap-mon-toan-lop-11-theo-chuong-trinh-giao-duc-pho-thong-2018.pdf (Lớp Toán Thầy Duy, Phạm Lê Duy, THPT Chu Văn An), Chương 3 Bài 1, Phần 3, Câu 15 (tr.346)"
%% kien_thuc: "Lớp 11 Bài 15 (giới hạn của dãy số)"
%% trang_thai: chinh-thuc
%% ghi_chu: "Đã duyệt 03/11/2026."
\begin{ex}
Biết rằng
\[\lim\left(\sqrt[3]{n^3+3}-\sqrt{n^2+2}\right)=\lim\dfrac a{\sqrt[3]{(n^3+3)^2}+n^2+n\cdot\sqrt[3]{n^3+3}}-\lim\dfrac b{n+\sqrt{n^2+2}}\]
với $a;b$ là các số nguyên dương. Tính $T=2a+3b$.
\shortans{12}
\loigiai{\sloppy $\sqrt[3]{n^3+3}-n=\dfrac{3}{\sqrt[3]{(n^3+3)^2}+n\sqrt[3]{n^3+3}+n^2}$ và $\sqrt{n^2+2}-n=\dfrac2{n+\sqrt{n^2+2}}$. Do đó $\sqrt[3]{n^3+3}-\sqrt{n^2+2}=\left(\sqrt[3]{n^3+3}-n\right)-\left(\sqrt{n^2+2}-n\right)$ nên $a=3$, $b=2$ và $T=2\cdot3+3\cdot2=12$.}
\end{ex}
